%\cleardoubleemptypage
\begin{abstract}
Reaktionsfähige Benutzeroberflächen sind in datenintensiven Echtzeit-Anwendungen wie Monitoring-Dashboards ein zentrales Qualitätsmerkmal. React~19 und Svelte~5 verfolgen dabei grundlegend unterschiedliche Architekturansätze: React~19 kombiniert ein Virtual-DOM-basiertes Laufzeitmodell mit einem Compiler zur automatischen Memoisierung, während Svelte~5 compilerbasierte feingranulare Reaktivität einsetzt. Bestehende Vergleichsstudien basieren jedoch überwiegend auf älteren Framework-Versionen und untersuchen vorwiegend isolierte Operationen; empirische Daten unter kombinierter Dauerlast aus steigender UI-Dichte und hoher Update-Frequenz fehlen bislang. Diese Arbeit untersucht, wie sich React~19 und Svelte~5 hinsichtlich Interaktionslatenz (INP), Scripting-Aufwand und Speicherverbrauch unter kombiniert steigender UI-Dichte und Update-Frequenz unterscheiden und ab welcher Last-Kombination \textit{Long Tasks} den Main Thread dauerhaft blockieren. Dazu wurden in einem automatisierten Benchmark (Playwright, Chrome DevTools Protocol) 1920 Messläufe über acht UI-Dichtestufen ($n \in \{100 \ldots 10\,000\}$) und vier Update-Frequenzen (10 bis 62\,Hz) durchgeführt. Als Metriken dienten INP, Scripting-Aufwand pro Tick, JS-Heap-Delta sowie Long-Task-Zähler. Die Ergebnisse zeigen, dass kein Framework uneingeschränkt überlegen ist. Bei hoher UI-Dichte liefert Svelte~5 in der Mehrzahl der Szenarien geringeren Scripting-Aufwand und niedrigere INP-Werte; zudem liegt Sveltes \textit{Tipping Point} bei höheren $n$-Werten. Über alle Dichtestufen weist Svelte~5 zudem einen 2,8- bis 6,4-fach geringeren Heap-Verbrauch auf. Allerdings erzielt React~19 bei moderater UI-Dichte und hoher Update-Frequenz durch kooperatives Yielding der Fiber-Engine niedrigere INP-Werte. Unter Extremlast ($n = 10\,000$, 62\,Hz) kehrt sich das Long-Task-Verhältnis um: Svelte erzeugt dort mehr \textit{Long Tasks} als React. Für die Praxis ergibt sich: Svelte~5 ist bei hoher UI-Dichte und in speicherkritischen Umgebungen vorzuziehen; React~19 kann bei moderater Dichte und hoher Update-Frequenz die bessere Wahl sein.
\end{abstract}

%\cleardoubleemptypage
\renewcommand\abstractname{Abstract}
\begin{abstract}
Responsive user interfaces are a central quality criterion in
data-intensive real-time applications such as monitoring dashboards.
React~19 and Svelte~5 pursue fundamentally different architectural
approaches: React~19 combines a virtual DOM-based runtime model with
a compiler for automatic memoization, while Svelte~5 employs
compiler-based fine-grained reactivity. However, existing comparison
studies are predominantly based on older framework versions and
examine primarily isolated operations; empirical data under combined
sustained load from increasing UI density and high update frequency
are lacking. This work investigates how React~19 and Svelte~5 differ
with respect to interaction latency (INP), scripting effort, and
memory consumption under combined increasing UI density and update
frequency, and at which load combination \textit{Long Tasks} persistently
block the main thread. To this end, 1920 measurement runs were
conducted in an automated benchmark (Playwright, Chrome DevTools
Protocol) across eight UI density levels ($n \in \{100 \ldots
10\,000\}$) and four update frequencies (10 to 62\,Hz). The metrics
used were INP, scripting effort per tick, JS heap delta, and long
task count. The results show that neither framework is universally
superior. At high UI density, Svelte~5 delivers lower scripting
effort and lower INP values in the majority of scenarios; Svelte's
\textit{tipping point} additionally lies at higher $n$ values. Across all
density levels, Svelte~5 additionally exhibits a 2.8- to 6.4-fold lower heap
consumption. However, React~19 achieves lower INP values at moderate
UI density and high update frequency through cooperative yielding of
the Fiber engine. Under extreme load ($n = 10\,000$, 62\,Hz), the long
task ratio reverses: Svelte produces more \textit{long tasks} than React. In practice: Svelte~5
is preferable at high UI density and in memory-constrained
environments; React~19 may be the better choice at moderate density
and high update frequency.
\end{abstract}